\chapter{Research Methodology}
\label{ch:methodology}
% How do music producers use stereo field?

% RQ1: How does the stereo image of a recording relate to real-
% world musical knowledge?

% RQ2: What methods can be used to describe the stereo image of
% a music recording?

In this chapter I break down the reasoning behind the \glspl{rq} given in \cref{sec:research_questions}, both from philosophical and practical standpoints. I argue that the methods I chose for the research in this thesis allow for a timely and achievable answer to the overarching \gls{rq}.

The overarching \gls{rq} -- \q{How do music producers use stereo field} -- could be approached in many ways. However, the work in this thesis is steeped in the context of large audio datasets. Indeed, the reason for performing the research at the current point in time is the ease of access to said datasets along with the availability of computer processing power to analyse them. From an ontological standpoint this most closely fits the Parmenidean viewpoint of \textit{being}. According to \citeauthor{grayDoingResearchReal2008}, Parmenidean ontology posits that reality is composed of clearly formed entities with identifiable properties \cite{grayDoingResearchReal2008}. If the producers' collective musical output is \q{reality}, then the recordings in the datasets certainly fit that description.

The two smaller \glspl{rq} represent different epistemological perspectives on the overarching one. \rqa{} is based in \textit{constructivism}, where the meaning of the stereo field is imposed on the recordings by the listener. Conversely, \rqb{} is rooted in \textit{objectivism}, where the properties of a recording's stereo field are an objective truth that can be measured. Not only that, but generalizations about the use of stereo field across the entire population of recordings are also objective and measurable. While the two epistemologies are somewhat adverse to each other, they both fall under the ontology of \textit{being} \cite{chiaProductionManagementKnowledge2002}.

Under \textit{constructivism}, one valid research methodology is \q{heuristic inquiry} -- the act of totally immersing oneself in the phenomenon being researched and allowing intuition to guide the understanding of the \gls{rq}. I used this technique for a sizeable part of the work described in this thesis. For example, {paper~\ref{paper:auxtrument}} is a retrospective account of building a new musical instrument where the primary sonic material was the signal from a remote stereo microphone. The act of building the instrument, listening to its output, and forming ideas about which sonic phenomena were interesting was a direct way to relate the stereo field with musical knowledge, addressing \rqa{}. While this provides a basis for other research, critics of heuristic inquiry will point out that it is deeply personal and therefore not particularly generalizable \cite{grayDoingResearchReal2008}.

To complement the \textit{constructive} approach and form more robust conclusions about an entire population of music producers, I concurrently embarked on a more objective research path. The epistemology of \textit{operationalism} is defined by the principle that in order to understand something, it must be measurable \cite{changOperationalism2021}. Thus, to understand stereo field, a set of operations to measure it is necessary. In practice, that means a software library capable of analysing music recordings and extracting information about their stereo field. The benefits of creating this software are threefold. Firstly, it enables comparative measurements of disparate music recordings, allowing one to validate and challenge the existing research described in \cref{sec:related_works}. Secondly, the methods used for any experimental research must be controllable and repeatable. A software library is a fantastic way to achieve this. Thirdly, the act of writing the library can be seen as a similar act of immersion to that previously described in the \textit{constructive} research. By breaking down the measurement of stereo field into simple instructions, one is forced to think about quantifiable relationships between audio phenomena. \citeauthor{wiltschkoLanguageThoughtCommunication2022} argues that language is a tool for both thought and communication in equal measure \cite{wiltschkoLanguageThoughtCommunication2022}. It follows that computer code (which is a written language like any other) acts as both a representation of the creator's thought process, and the communication of those thoughts to any readers. The software written for this thesis was therefore fundamental to the research process, as well as underpinning the results of paper \ref{paper:dafx}.

This combinatorial approach is not a new idea. The 13th century philosopher Nicolaus Cusanus spoke about two foundations of knowledge - \q{ratio} (that which you measure) and \q{intellectus} (the relationship to the unknown) \cite{millerCusanusNicolausNicolas2025}.  It captures the dichotomy that when measuring a phenomenon, more measurements will give one a more accurate picture of the truth, but one must have the forethought to imagine that the measurements might not be able to represent the entire truth. Cusanus argued that both are necessary to make sense of the world, and modern academics propose that as our world is increasingly measured and rationalized, intellectus becomes increasingly important \cite{hylvingRatioIntellectusConceptual2022}.

One potential criticism of this combinatorial approach is that an answer to the overarching \gls{rq} could instead be gained by gathering interview or observational data (using \citeauthor{yinCaseStudyResearch2014}'s taxonomy of \q{evidence sources} \cite{yinCaseStudyResearch2014}).  As \citeauthor{grayDoingResearchReal2008}'s assessment of the strengths and weaknesses of each evidence type states \cite{grayDoingResearchReal2008}, the former (interviews and observation) result in deep, contextually rich data, but are time-consuming to collect and subject to biases from the act of interviewing or observing itself. More specific to the music community, there is an inherent understanding that music is a form of language \cite{agawuMusicDiscourseSemiotic2014}. Adding a layer of indirection (interviews conducted in spoken language) only serves to muddy the resulting data. Indeed, \citeauthor{hospersConceptArtisticExpression1954} argues that an artist's expression is inherently intertwined with the creation the original art form, noting that the intentions are sometimes so complex and multifaceted as to defy explanation even to the artist themselves \cite{hospersConceptArtisticExpression1954}. In contrast, the musical datasets that we are presented with would be classified as \q{archival records} or \q{physical artefacts}, which according to \citeauthor{grayDoingResearchReal2008} are precise and technically insightful, but prone to selectivity and reporting biases \cite{grayDoingResearchReal2008}. 

I argue in this thesis that such biases can be partially mitigated by combining research from multiple epistemologies. While \textit{objective} research can determine what music producers do from a technical perspective, \textit{constructive} research can offer insight into the choices and trade-offs they made to end up at their final artistic output. The work I have performed for this thesis is a combination of immersion in musical practice, thorough reviews of existing literature, and detailed implementation of audio \gls{dsp} analysis techniques. This combination presents a well-rounded viewpoint on the original \glspl{rq}, as well as giving rise to new questions that could be answered in the continuation to doctoral studies (see \cref{sec:future_work}).
